% ==============================================================================
\chapter{Estado del arte}
\label{ch:estado-del-arte}
% ==============================================================================

Este capítulo revisa aplicaciones cívicas relevantes para EcoAlerta, compara las tecnologías candidatas para cada capa y justifica las decisiones tomadas en el diseño del Capítulo~\ref{ch:analisis-disenyo}.


\section{Contexto: aplicaciones cívicas geolocalizadas}
\label{sec:contexto-apps}

Las aplicaciones cívicas con reporte geolocalizado son parte del movimiento \textit{civic tech} y de lo que la literatura llama \textit{ciudadanos como sensores} \parencite{goodchild2007citizens, haklay2013citizen}. Tres innovaciones las hicieron viables a partir de mediados de los 2000:

\begin{enumerate}
  \item \textit{Smartphones} con cámara y GPS (iPhone 2007, Android 2008): cualquier ciudadano puede generar evidencia fotográfica geolocalizada en tiempo real.
  \item OpenStreetMap (2004): cartografía libre, sin dependencia de proveedores comerciales.
  \item Datos abiertos y leyes de transparencia (en España, Ley 19/2013): marco para el intercambio bidireccional ciudadano--administración.
\end{enumerate}

A partir de ahí se consolidan dos familias: las \textbf{generalistas} (FixMyStreet, SeeClickFix, que cubren cualquier incidencia urbana) y las \textbf{especializadas} (iNaturalist para biodiversidad, Línea Verde para reporte municipal en España). EcoAlerta se sitúa en un hueco intermedio: especializada en medio ambiente, con componente social.


\section{Aplicaciones analizadas}
\label{sec:apps-analizadas}

Se han seleccionado cinco aplicaciones por relevancia, heterogeneidad y cobertura del espectro de la \textit{civic tech} aplicada al medio ambiente: FixMyStreet, Línea Verde, iNaturalist, SeeClickFix y Epicollect5. El análisis de cada una se ha realizado a partir de la documentación oficial, las páginas web del producto y reseñas públicas de usuarios, complementadas con la observación de capturas de pantalla y vídeos demostrativos disponibles en línea. A continuación se describe cada una de ellas.


\subsection{FixMyStreet}
\label{subsec:fixmystreet}

Desarrollada en 2007 por \textit{mySociety}~\parencite{fixmystreet}, organización británica sin ánimo de lucro especializada en \textit{civic tech}. El código es abierto bajo licencia AGPL-3.0 y la plataforma ha sido desplegada en cientos de ayuntamientos del Reino Unido, así como en versiones locales para Bruselas, Noruega o Suiza.

\textbf{Funcionalidades principales.}  Reporte ciudadano de incidencias urbanas (baches, alumbrado, mobiliario dañado, vertidos) con foto opcional y ubicación geoespacial, listado público de reportes por zona, panel administrativo para los ayuntamientos contratantes, notificaciones de cambio de estado al autor del reporte y duplicado detectado automáticamente por proximidad.

\textbf{Flujo de uso.}  El usuario abre la web móvil, introduce una dirección o selecciona un punto en el mapa, escoge una categoría y describe la incidencia. La plataforma encamina el reporte al ayuntamiento competente según la ubicación del marcador y le asigna un identificador público de seguimiento.

\textbf{Tecnología.}  Aplicación web con diseño \textit{mobile-friendly}. \textit{Backend} en Perl con Plack y PostgreSQL. Cartografía sobre OpenStreetMap.

\textbf{Aportaciones de referencia para EcoAlerta.}  Modelo de delegación municipal por ubicación geoespacial, flujo de estados visible para el ciudadano, identificador público de seguimiento y comunidad de usuarios.

\textbf{Limitaciones respecto al objetivo de EcoAlerta.}  Enfoque generalista (urbano) sin especialización medioambiental. No contempla entidades responsables distintas del ayuntamiento (SEPRONA, bomberos, ONG). No incluye severidad con impacto en la priorización ni componente social significativo (votos, seguidores).


\subsection{Línea Verde}
\label{subsec:linea-verde}

Servicio comercial de la empresa española \textit{Línea Verde Configuración y Servicios Medioambientales S.A.}~\parencite{lineaverde}, contratado por más de 300 municipios españoles. Funciona bajo modelo B2G \textit{white-label}: cada Ayuntamiento adquiere el servicio y obtiene su propia versión personalizada con marca municipal.

\textbf{Funcionalidades principales.}  Reporte de quejas e incidencias municipales con foto y ubicación, consulta del calendario de recogida de residuos, búsqueda de puntos de reciclaje, alertas sobre cortes de servicios y avisos del ayuntamiento. La rama medioambiental cubre vertidos, animales muertos en vía pública, mobiliario urbano dañado y problemas de jardinería.

\textbf{Flujo de uso.}  Aplicación móvil nativa (iOS/Android) y portal web. El ciudadano selecciona el tipo de incidencia, describe el problema, adjunta foto y envía. El reporte queda en cola del ayuntamiento contratante, que lo gestiona desde su propio panel y notifica al usuario cuando cambia el estado.

\textbf{Tecnología.}  Aplicación móvil nativa publicada en App Store y Google Play, con portal web complementario. \textit{Backend} y arquitectura no documentados públicamente al tratarse de software propietario.

\textbf{Aportaciones de referencia para EcoAlerta.}  Modelo \textit{white-label} viable para administraciones locales españolas. Estructura de categorización compatible con la legislación municipal española. Notificaciones de cambio de estado al ciudadano.

\textbf{Limitaciones respecto al objetivo de EcoAlerta.}  Software propietario sin código fuente público; cada despliegue queda aislado en el ámbito municipal contratante; el reporte no se delega a entidades especializadas distintas del propio ayuntamiento; ausencia de componente social entre ciudadanos.


\subsection{iNaturalist}
\label{subsec:inaturalist}

Plataforma global de ciencia ciudadana para observación de biodiversidad nacida en 2008 en la Universidad de California~\parencite{inaturalist}. Actualmente está gestionada conjuntamente por la \textit{California Academy of Sciences} y la \textit{National Geographic Society}. Acumula más de 180 millones de observaciones aportadas por una comunidad internacional de naturalistas y voluntarios.

\textbf{Funcionalidades principales.}  Subida de observaciones de fauna y flora con foto y geolocalización, identificación automática de especies mediante una red neuronal entrenada con el corpus comunitario, validación comunitaria de las identificaciones, proyectos colaborativos por región o ecosistema y exportación de datos para investigación científica.

\textbf{Flujo de uso.}  El usuario abre la aplicación, fotografía un ejemplar, la app sugiere especies candidatas, el usuario confirma o ajusta y publica. Otros usuarios pueden validar la identificación, lo que eleva el grado de confianza de la observación hasta «calidad de investigación».

\textbf{Tecnología.}  Aplicación móvil nativa y web. \textit{Backend} en Ruby on Rails. Base de datos PostgreSQL con extensión PostGIS para datos geoespaciales. Elasticsearch para búsquedas. Servicio de visión por computador propio para reconocimiento de imagen.

\textbf{Aportaciones de referencia para EcoAlerta.}  Comunidad activa que valida los reportes, perfil público del usuario con histórico de observaciones, sistema de proyectos por región o temática, uso de PostGIS para consultas geoespaciales (decisión arquitectónica equivalente a la adoptada en EcoAlerta).

\textbf{Limitaciones respecto al objetivo de EcoAlerta.}  Enfoque en biodiversidad, no en incidencias medioambientales. No contempla delegación a entidades responsables ni flujo de estados administrativo. No prioriza por severidad ni por urgencia de actuación.


\subsection{SeeClickFix}
\label{subsec:seeclickfix}

Fundado en 2008 en New Haven (EE.UU.), referente del reporte cívico americano~\parencite{seeclickfix}. Adquirido por CivicPlus en 2019, sigue operando en municipios norteamericanos.

\textbf{Funcionalidades.}  Reporte geolocalizado con foto, categorías estandarizadas, flujo de estados (recibido, reconocido, en curso, resuelto), panel administrativo, integración con sistemas \textit{backend} municipales vía API.

\textbf{Tecnología.}  App móvil nativa (iOS/Android) y web. API documentada para integraciones.

\textbf{Relevancia para EcoAlerta.}  Valida el modelo B2G con flujo de estados trazable. Su limitación es el foco urbano exclusivo: no aborda incendios forestales, fauna silvestre ni daños a espacios naturales.


\subsection{Epicollect5}
\label{subsec:epicollect}

Proyecto del Imperial College London para recolección de datos de campo en investigación científica~\parencite{epicollect5}. Usado en epidemiología, conservación y ecología.

\textbf{Funcionalidades.}  Formularios personalizados, recolección \textit{offline} con sincronización diferida, exportación en CSV/JSON. Gratuito para investigación.

\textbf{Tecnología.}  App móvil y web. \textit{Stack} no documentado públicamente.

\textbf{Relevancia para EcoAlerta.}  Su robustez \textit{offline}-\textit{first} es referencia para el RF-19. Su enfoque genérico lo hace demasiado amplio para el ciudadano medio, pero útil para ONG conservacionistas que pueden configurar formularios específicos.


\section{Comparativa funcional}
\label{sec:comparativa-funcional}

\begin{longtable}{@{}L{4.5cm}C{1.4cm}C{1.4cm}C{1.4cm}C{1.4cm}C{1.4cm}C{1.4cm}@{}}
\caption{Comparativa funcional. $\checkmark${} sí, $\sim${} parcial, $\times${} no.} \label{tab:comparativa-funcional} \\
\toprule
\textbf{Funcionalidad} & \textbf{FixMy\-Street} & \textbf{Línea Verde} & \textbf{iNatu\-ralist} & \textbf{SeeClick\-Fix} & \textbf{Epi\-collect} & \textbf{Eco\-Alerta} \\
\midrule
\endfirsthead
\toprule
\textbf{Funcionalidad} & \textbf{FixMy\-Street} & \textbf{Línea Verde} & \textbf{iNatu\-ralist} & \textbf{SeeClick\-Fix} & \textbf{Epi\-collect} & \textbf{Eco\-Alerta} \\
\midrule
\endhead
\bottomrule
\endfoot
Reporte geolocalizado con foto    & $\checkmark$ & $\checkmark$ & $\checkmark$ & $\checkmark$ & $\checkmark$ & $\checkmark$ \\
Enfoque medioambiental específico & $\times$ & $\sim$ & $\checkmark$ & $\times$ & $\sim$ & $\checkmark$ \\
Delegación a entidades            & $\sim$ & $\checkmark$ & $\times$ & $\checkmark$ & $\times$ & $\checkmark$ \\
Flujo de estados trazable         & $\checkmark$ & $\sim$ & $\times$ & $\checkmark$ & $\times$ & $\checkmark$ \\
Componente social (votos, comentarios) & $\sim$ & $\times$ & $\checkmark$ & $\sim$ & $\times$ & $\checkmark$ \\
Perfil público de usuario         & $\times$ & $\times$ & $\checkmark$ & $\sim$ & $\times$ & $\checkmark$ \\
Acceso anónimo de consulta        & $\checkmark$ & $\checkmark$ & $\checkmark$ & $\checkmark$ & $\sim$ & $\checkmark$ \\
Niveles de severidad explícitos   & $\sim$ & $\times$ & $\times$ & $\sim$ & $\checkmark$ & $\checkmark$ \\
\textit{Audit trail} completo     & $\times$ & $\times$ & $\times$ & $\sim$ & $\times$ & $\checkmark$ \\
PWA instalable                    & $\sim$ & $\times$ & $\times$ & $\times$ & $\sim$ & $\checkmark$ \\
\textit{Offline}                  & $\times$ & $\times$ & $\checkmark$ & $\times$ & $\checkmark$ & $\checkmark$ \\
2FA / \textit{audit} de seguridad & $\times$ & $\times$ & $\times$ & $\times$ & $\times$ & $\checkmark$ \\
\textit{Open source}              & $\checkmark$ & $\times$ & $\checkmark$ & $\times$ & $\sim$ & $\checkmark$ \\
\end{longtable}

Ninguna aplicación cubre simultáneamente las seis dimensiones que busca EcoAlerta: enfoque medioambiental, delegación a entidades heterogéneas (no solo municipales), flujo trazable, social significativo, severidad con impacto en priorización y \textit{audit trail}. Ese hueco es el punto de partida del proyecto (objetivo 1).


\section{Tecnologías evaluadas}
\label{sec:comparativa-tecnica}

Para cada capa consideré varias opciones. La Tabla~\ref{tab:tecnologias} resume las decisiones y por qué.

\begin{table}[H]
\centering
\caption{Decisiones tecnológicas por capa.}
\label{tab:tecnologias}
\small
\begin{tabular}{@{}L{2.8cm}L{4cm}L{6cm}@{}}
\toprule
\textbf{Capa} & \textbf{Elección vs alternativas} & \textbf{Razón} \\
\midrule
Plataforma & PWA \textit{vs} app nativa & PWA con una sola \textit{code base}, distribución por URL, instalable. La nativa quedaría como trabajo futuro. \\
\textit{Frontend} & React \textit{vs} Vue, Angular & React tiene la mejor integración con Leaflet (\texttt{react-leaflet}), gran mercado laboral en España, y ya conocía la herramienta. \\
BD & PostgreSQL+PostGIS \textit{vs} MongoDB & PostGIS aporta \texttt{ST\_DWithin}, índices GIST y transacciones. El modelo relacional encaja mejor con las relaciones del dominio. MongoDB obligaría a denormalizar. \\
Mapa & Leaflet+OSM \textit{vs} Google Maps & Leaflet es \textit{open source}, sin coste ni dependencia comercial. Coherente con la filosofía \textit{civic tech} de las apps de referencia. \\
Metodología & SDD \textit{vs} Scrum clásico & Para un solo desarrollador con asistencia IA, la especificación previa funciona como guía al agente y como base para revisar. Más útil que un Scrum tradicional pensado para equipos. \\
\bottomrule
\end{tabular}
\end{table}


\section{Posicionamiento de EcoAlerta}
\label{sec:posicionamiento}

Una vez analizadas todas ellas, hay que tener en cuenta que EcoAlerta combina elementos de varias fuentes:

\begin{itemize}
  \item Flujo de estados trazable de FixMyStreet y SeeClickFix, ampliado con validación, marcado de duplicados y escalado por \textit{timeout}.
  \item Delegación B2G multi-entidad de Línea Verde, generalizada a entidades no municipales (SEPRONA, bomberos, ONG).
  \item Componente social de iNaturalist (seguimiento, votos, comentarios), adaptado al ámbito medioambiental con la distinción de comentarios oficiales.
  \item \textit{Offline} robusto al estilo Epicollect5, con Service Worker e IndexedDB.
  \item Mecanismos de seguridad propios (2FA TOTP, \textit{audit log}, Cloudflare Turnstile~\parencite{turnstile}) ausentes en las cinco apps analizadas.
\end{itemize}

Con eso, más la prioridad dada a la accesibilidad WCAG y a la transparencia (código abierto y declaración explícita del uso de IA), se cubre el hueco que las aplicaciones existentes dejan.
